\subsection{Đề 2}
\setcounter{ex}{0}
\Opensolutionfile{ans}[ans/ans-0-B8]
\caulc
%câu 1
\begin{ex}%[2H2N2-3]
	Trong không gian $Oxyz$, cho hai véctơ là $\vec{a}=(1;2;3)$ và $\vec{b}=(4;5;6)$. Toạ độ của véctơ $\vec{a}+\vec{b}$ tương ứng là
	\choice
	{\True $(5;7;9)$}
	{$(3;7;9)$}
	{$(5;3;9)$}
	{$(3;5;9)$}
	\loigiai{
		Ta có $\vec{a}+\vec{b}=(1+4;2+5;3+6)=(5;7;9)$.
	}
\end{ex}	
%câu 2
\begin{ex}%[2H2N2-3]
	Trong không gian $Oxyz$, cho hai véctơ là $\vec{a}=(2;-1;4)$. Tọa độ của  véctơ $-\vec{a}$ là
	\choice
	{\True $(-2;1;-4)$}
	{$(2;-1;-4)$}
	{$(-2;-1;4)$}
	{$(2;1;4)$}
	\loigiai{
		Ta có $-\vec{a}=(-2;1;-4)$.
	}
\end{ex}	
%câu 3
\begin{ex}%[2H2N2-4]
	Trong không gian $Oxyz$, cho hai véctơ là $\vec{u}=(1;2;3)$ và $\vec{v}=(4;-5;6)$. Tìm tích vô hướng của hai véctơ  $\vec{u}$ và $\vec{v}$.
	\choice
	{\True $12$}
	{$20$}
	{$14$}
	{$28$}
	\loigiai{
		Ta có $\vec{u}\cdot\vec{v}=1\cdot 4+2\cdot (-5)+3\cdot 6=12$.
	}
\end{ex}	
%câu 4
\begin{ex}%[2H2N2-2]
	Trong không gian $Oxyz$, cho điểm $A(1;2;3)$ và điểm $B(3;4;5)$. Tọa độ điểm $C$ là trung điểm của đoạn thẳng $AB$ là
	\choice
	{\True $(2;3;4)$}
	{$(1; 3;5)$}
	{$(2;4;3)$}
	{$(2;2;3)$}
	\loigiai{
		Ta có $\heva{&x_C=\dfrac{x_A+x_B}{2}=2\\&y_C=\dfrac{y_A+y_B}{2}=3\\&z_C=\dfrac{z_A+z_B}{2}=4.} $\\
		Vậy $C(2;3;4)$.
	}
\end{ex}	
%câu 5
\begin{ex}%[2H2N2-2]
	Trong không gian với hệ tọa độ $Oxyz$, cho ba điểm $A(3;2;-5)$, $B(1;2;4)$, $C(2;5;-2)$. Tọa độ trọng tâm $G$ của $\triangle ABC$ là
	\choice
	{$G(6;9;-3)$} 
	{\True $G(2;3;-1)$}
	{$G(2;3;1)$} 
	{$G(6;9;3)$}
	\loigiai
	{
		Ta có $\heva{& x_G=\dfrac{3+1+2}{3}\\&y_G=\dfrac{2+2+5}{3}\\&z_G=\dfrac{-5+4-2}{3} }\Leftrightarrow\heva{& x_G=2\\&y_G=3\\&z_G=-1 .}$\\
		Vậy $G(2;3;-1)$.
	}
\end{ex}

%câu 6
\begin{ex}%[2H2H2-3]
	Trong không gian $Oxyz$, cho hai véctơ là $\vec{a}=(1;2;3)$ và $\vec{b}=(2;-1;1)$. Tìm biểu thức tọa độ của véctơ  $\vec{c}=3\vec{a}-2\vec{b}$.
	\choice
	{\True $(-1;8;7) $}
	{$(3;4;5)$}
	{$(-1;7;7)$}
	{$(4;5;6)$}
	\loigiai{
		Ta có $\vec{c}=3\vec{a}-2\vec{b}=\left(3\cdot 1-2\cdot 2; 3\cdot 2-2\cdot (-1);3\cdot 3-2\cdot 1\right)=(-1;8;7)$.
	}
\end{ex}	 
%câu 7
\begin{ex}%[2H2H2-2]
Trong không gian $Oxyz$, điểm đối xứng với điểm $M(2;2;-1)$ qua mặt phẳng $(Oyz)$ là
\choice
{$M'(2;-2;1)$}
{$M'(-2;0;0)$}
{$M'(-2;-2;1)$}
{\True $M'(-2;2;-1)$}
\loigiai{
	Điểm đối xứng với điểm $M(2;2;-1)$ qua mặt phẳng $(Oyz)$ là $M'(-2;2;-1)$.
}
\end{ex}

%câu 8
\begin{ex}%[2H2H2-2]
	Trong không gian $Oxyz$, cho điểm $M(1;-2;3)$. Điểm đối xứng với $M$ qua trục $Oy$ là
	\choice
	{$M_{2}(1;2;3)$}
	{$M_{3}(-1;2;-3)$}
	{\True $M_{4}(-1;-2;-3)$}
	{$M_{1}(1;-2;-3)$}
	\loigiai{
		Ta có $M(1;-2;3)$.\\
		Vậy điểm đối xứng với $M$ qua qua trục $Oy$ là $M_{4}(-1;-2;-3)$.
	}
\end{ex}

%câu 9
\begin{ex}%[2H2H2-2]
	Trong không gian với hệ trục tọa độ $Oxyz$ cho tam giác $ABC$, với $A(1;2;1)$, $B(-3;0;3)$, $C(2;4;-1)$. Tìm tọa độ điểm $D$ sao cho tứ giác $ABCD$ là hình bình hành.		
	\choice
	{\True $D(6;6;-3)$}
	{$D(6;-6;3)$}
	{$D(6;6;3)$}
	{$D(6;-6;-3)$}
	\loigiai{
		Gọi $D(x;y;z)$.\\
		Ta có $\overrightarrow {AB}=(-4;-2;2)$; $\overrightarrow {DC}=(2-x;4-y;-1-z)$.\\
		$ABCD$	 là hình bình hành $\Leftrightarrow \overrightarrow {AB}=\overrightarrow {DC}\Leftrightarrow\heva{&-4=2-x\\&-2=4-y\\&2=-1-z}\Leftrightarrow\heva{&x=6\\&y=6\\&z=-3.}$\\
		Vậy $D(6;6;-3)$.
	}
\end{ex}

%câu 10
\begin{ex}%[2H2H2-4]
	Trong không gian $Oxyz$, cho hai véctơ là $\vec{a}=(1;-2;3)$ và $\vec{b}=(-2;1;2)$. Tích vô hướng $\left(\vec{a}+\vec{b}\right)\cdot \vec{b}$ bằng
	\choice
	{$12$}
	{$2$}
	{\True $11$}
	{$10$}
	\loigiai{
	Ta có 	$\vec{a}+\vec{b}=(-1;-1;5)$.\\
	Suy ra $\left(\vec{a}+\vec{b}\right)\cdot \vec{b}=(-1)\cdot (-2)+(-1)\cdot 1+5\cdot 2=11$.
	}
\end{ex}	
%câu 11
\begin{ex}%[2H2H2-4]
	Trong không gian $Oxyz$, cho hai véctơ là $\vec{a}=(1;0;0)$ và $\vec{b}=\left(\dfrac{1}{2};\dfrac{\sqrt{3}}{2};0\right)$. Góc giữa hai véctơ này là
	\choice
	{$30^\circ$}
	{$45^\circ$}
	{\True $60^\circ$}
	{$90^\circ$}
	\loigiai{
	Ta có
	\allowdisplaybreaks
	\begin{eqnarray*}
		\cos\left(\overrightarrow{a},\overrightarrow{b}\right)
		&=&\dfrac{\overrightarrow{a}\cdot \overrightarrow{b}}{\left|\overrightarrow{a}\right|\cdot \left|\overrightarrow{b}\right|}
		\\
		&=&\dfrac{1\cdot \frac{1}{2}+0\cdot \frac{\sqrt{3}}{2}+0\cdot 0}{\sqrt{1^2+0^2+0^2}\cdot \sqrt{\left(\frac{1}{2}\right)^2+\left(\frac{\sqrt{3}}{2}\right)^2+0^2}}
		\\
		&=&\dfrac{1}{2}.
	\end{eqnarray*}
	Suy ra $\left(\vec{a},\vec{b}\right)=60^\circ$.	
	}
\end{ex}	
% Câu 12
\begin{ex}%[2H2V2-4]
	Trong không gian $Oxyz$, cho hai điểm $A(1;2;1)$; $B(2;-1;3)$ và điểm $M(a;b;0)$ sao cho $MA^2+MB^2$ nhỏ nhất. Giá trị của $a+b$ là
	\choice
	{$3$}
	{$-2$}
	{$1$}
	{\True $2$}
	\loigiai{
		Ta thấy $M(a;b;0)\in(Oxy)$.\\
		Gọi $I\left(\dfrac{3}{2};\dfrac{1}{2};2\right)$ là trung điểm của đoạn thẳng $AB$, ta có
		\begin{eqnarray*}
			MA^2+MB^2&=&\overrightarrow{MA} ^2+\overrightarrow{MB} ^2=\left(\overrightarrow{IA} -\overrightarrow{IM}\right)^2+\left(\overrightarrow{IB} -\overrightarrow{IM}\right)^2\\
			&=&\left(\overrightarrow{IA} ^2+\overrightarrow{IM} ^2-2\overrightarrow{IA} \cdot\overrightarrow{IM}\right)+\left(\overrightarrow{IB} ^2+\overrightarrow{IM} ^2-2\overrightarrow{IB} \cdot\overrightarrow{IM}\right)\\
			&=&2IM^2+2IA^2=2IM^2+7.
		\end{eqnarray*}
		Suy ra $MA^2+MB^2$ nhỏ nhất $\Leftrightarrow $ $IM$ ngắn nhất $\Leftrightarrow $ $M$ là hình chiếu vuông góc của $I$ trên mặt phẳng $(Oxy)$. \\
		Suy ra $M\left(\dfrac{3}{2};\dfrac{1}{2};0\right)$. \\
		Như vậy $a=\dfrac{3}{2},b=\dfrac{1}{2}\Rightarrow a+b=\dfrac{3}{2}+\dfrac{1}{2}=2$.
	}
\end{ex}

\Closesolutionfile{ans}
\indapan{6}{ans/ans-0-B8}
\cauds
\Opensolutionfile{ans}[ans/ans-0-B8-DS]
\setcounter{ex}{0}
%Câu 1
\begin{ex}%[2H2V2-2]
	Trong không gian $Oxyz$, cho các điểm  $A(1;-2;3)$, $B(-2;1;2)$, $C(3;-1;2)$.
	\choiceTF
	{$\overrightarrow{AB}=(-3;3;1)$}
	{$\overrightarrow{AB}=3\overrightarrow{AC}$}
	{\True Ba điểm $A$, $B$, $C$ không thẳng hàng}
	{Tọa độ chân đường cao vẽ từ $A$ của tam giác $ABC$ là $\left(-\dfrac{47}{29};\dfrac{13}{29};2\right)$}
	\loigiai{
		\begin{itemchoice}
			\itemch Sai. Ta có $\overrightarrow{AB}=(-3;3;-1)$.
			\itemch Sai. Ta có $\overrightarrow{AC}=(2;1;-1)\Rightarrow 3\overrightarrow{AC}=(6;3;-3)$. Do đó $\overrightarrow{AB}\neq 3\overrightarrow{AC}$.
			\itemch Đúng. Vì $\overrightarrow{AB}=(-3;3;1)$, $\overrightarrow{AC}=(2;1;-1)$ suy ra $\overrightarrow{AB}$, $\overrightarrow{AC}$ không cùng phương. \\
			Vậy $A$, $B$, $C$ không thẳng hàng.
			\itemch Sai. Gọi $A'(x;y;z)$ là chân đường cao vẽ từ $A$ của tam giác $ABC$.\\
			Ta có $\heva{&\overrightarrow{AA'}\cdot \overrightarrow{BC}=0&\qquad(1)\\&\overrightarrow{BC}, \overrightarrow{BA'} ~\text{cùng phương}.&\qquad(2)} $\\
			Ta có $\overrightarrow{AA'}=(x-1;y+2;z-3)$, $\overrightarrow{BC}=(5;-2;0)$, $\overrightarrow{BA'}=(x+2;y-1;z-2)$.\\
			$(1)\Leftrightarrow 5\cdot (x-1)-2\cdot (y+2)+0\cdot (z-3)=0 \Leftrightarrow 5x-2y=9$.\\
			$(2)\Leftrightarrow\heva{&\dfrac{x+2}{5}=\dfrac{y-1}{-2}\\&z-2=0}\Leftrightarrow\heva{&2x+5y=1\\&z=2.}$\\
			Suy ra $\heva{&x=\dfrac{47}{29}\\&y=-\dfrac{13}{29}\\&z=2}\Rightarrow A'\left(\dfrac{47}{29};-\dfrac{13}{29};2\right)$.
		\end{itemchoice}
	}
\end{ex}
%Câu 2
\begin{ex}%[2H2V2-2]
	Trong không gian $Oxyz$, cho các điểm  $A(1;3;5)$, $B(1;1;3)$, $C(4;-2;3)$.
	\choiceTF
	{\True Tọa độ trung điểm của $BC$ là $\left(\dfrac{5}{2};-\dfrac{1}{2};3\right) $}
	{\True Độ dài đoạn thẳng $BC$ là $3\sqrt{2}$}
	{\True Côsin $\widehat{BAC}$ bằng $\dfrac{7\sqrt{19}}{38}$}
	{Gọi $D$ là đỉnh thứ tư của hình bình hành $ABCD$. Tọa độ hình chiếu của trọng tâm tam giác $ABD$ lên mặt phẳng $Oyz$ là $(2;0;0)$}
	\loigiai{
		\begin{itemchoice}
			\itemch Đúng.  Gọi $M$ là trung điểm $BC$. \\
			Ta có $\heva{&x_M=\dfrac{x_B+x_C}{2}=\dfrac{1+4}{2}=\dfrac{5}{2}\\&y_M=\dfrac{y_B+y_C}{2}=\dfrac{1-2}{2}=-\dfrac{1}{2}\\&z_M=\dfrac{z_B+z_C}{2}=\dfrac{3+3}{2}=3}$ suy ra $M\left(\dfrac{5}{2};-\dfrac{1}{2};3\right) $.
			\itemch Đúng. Ta có $BC=\sqrt{(4-1)^2+(-2-1)^2+(3-3)^2}=3\sqrt{2}$.
			\itemch Sai. Ta có $\overrightarrow{AB}=(0;-2;-2)$, $\overrightarrow{AC}=(3;-5;-2)$.
			\allowdisplaybreaks
			\begin{eqnarray*}
			 \cos  \widehat{BAC}&=&\cos\left(\overrightarrow{AB},\overrightarrow{AC}\right)
			 \\
			 &=&\dfrac{\overrightarrow{AB}\cdot \overrightarrow{AC}}{AB\cdot AC}
			 \\
			 &=&\dfrac{0\cdot 3+(-2)\cdot (-5)+(-2)\cdot (-2)}{\sqrt{0^2+(-2)^2+(-2)^2}\cdot\sqrt{3^2+(-5)^2+(-2)^2}}
			 \\
			 &=&\dfrac{7\sqrt{19}}{38}.			 
		 \end{eqnarray*}	 
			\itemch Sai. Gọi $D(x;y;z)$.\\
			Ta có $\overrightarrow {AB}=(0;-2;-2)$; $\overrightarrow {DC}=(4-x;-2-y;3-z)$.\\
			$ABCD$	 là hình bình hành $\Leftrightarrow \overrightarrow {AB}=\overrightarrow {DC}\Leftrightarrow\heva{&0=4-x\\&-2=-2-y\\&-2=3-z}\Leftrightarrow\heva{&x=4\\&y=0\\&z=5.}$\\
			Suy ra $D(4;0;5)$.\\
			Gọi $G$ là trọng tâm tam giác $ABD$.\\
			Ta có $\heva{& x_G=\dfrac{1+1+4}{3}\\&y_G=\dfrac{3+1+0}{3}\\&z_G=\dfrac{5+3-1}{3} }\Leftrightarrow\heva{& x_G=2\\&y_G=\dfrac{4}{3}\\&z_G=\dfrac{7}{3}.}$\\
			Suy ra $G\left(2;\dfrac{4}{3};\dfrac{4}{3}\right)$.\\
			Vậy hình chiếu  của $G$  lên mặt phẳng $Oyz$ là $H\left(0;\dfrac{4}{3};\dfrac{4}{3}\right)$.
		\end{itemchoice}
	}
\end{ex}
%Câu 3
\begin{ex}%[2H2V2-2]
	Trong không gian với hệ trục tọa độ $Oxyz$, cho $\overrightarrow{OA}=3\overrightarrow{i}-\overrightarrow{k}$, với $\overrightarrow{i}$, $\overrightarrow{k}$  là hai véctơ đơn vị trên hai trục tọa độ $Ox$, $Oz$, hai điểm  $B(-1;2;3)$, $C(1;4;1)$. 	
	\choiceTF
	{\True $A(3;0;-1)$}
	{Ba điểm $A$, $B$, $C$ thẳng hàng}
	{\True Điểm $D(a;b;c)$ là điểm đối xứng của với $A$ qua $B$. Khi đó $a+b+c=6$}
	{\True Điểm $M(m;n;p)$ trên mặt phẳng $(Oxy)$ sao cho $MA^2+MB^2+MC^2$ đạt giá trị nhỏ nhất. Khi đó $2m-n+2024p=0$}
	\loigiai{
		\begin{itemchoice}
			\itemch Đúng. Vì $\overrightarrow{OA}=3\overrightarrow{i}-\overrightarrow{k}$ nên $A(3;0;-1)$.
			\itemch Sai. Ta có $\overrightarrow{AB}=(4;2;4)$, $\overrightarrow{AC}=(-2;4;2)$.\\
			Vì $4:2:4\neq -2:4:2$ nên $\overrightarrow{AB}$, $\overrightarrow{AC}$ không cùng phương. Suy ra $A$, $B$, $C$ không thẳng hàng.
			\itemch Đúng. Vì $D$ là điểm đối xứng với $A$ qua $B$ nên  $B$ là trung điểm của $AD$.\\ Ta có $\heva{& x_D=2x_B-x_A=-5 \\ & y_D=2y_B-y_A=4 \\ & z_D=2z_B-z_A=7.}$\\
			Suy ra $D(-5;4;7)$. Do đó $a=-5$, $b=4$, $c=7$.\\
			Vậy $a+b+c=6$.
			\itemch Đúng. Gọi $I(x;y;z)$ là điểm thỏa mãn $\overrightarrow{IA} + \overrightarrow{IB} + \overrightarrow{IC} = \overrightarrow{0}$. Ta có 
			$$\heva{&3-x-1-x+1-x=0\\&0-y+2-y+4-y=0\\&-1-z+3-z+1-z=0} \Leftrightarrow \heva{&x=1\\&y=2\\&z=1} \Rightarrow I(1;2;1).$$	
			\begin{eqnarray*}
				MA^2+MB^2+MC^2 &=& \left(\overrightarrow{MI} + \overrightarrow{IA}	\right)^2 + \left(\overrightarrow{MI} + \overrightarrow{IB}	\right)^2 +\left(\overrightarrow{MI} + \overrightarrow{IC}	\right)^2 \\
				&=& 3MI^2+IA^2+IB^2+IC^2 + 2\overrightarrow{MI} \left(\overrightarrow{IA} + \overrightarrow{IB} + \overrightarrow{IC} \right)\\
				&=& 3MI^2+IA^2+IB^2+IC^2.
			\end{eqnarray*}
			Do $IA^2+IB^2+IC^2$ không thay đổi nên $ MA^2+MB^2+MC^2$ nhỏ nhất khi $MI$ nhỏ nhất. Hay $M$ là hình chiếu của điểm $I$ trên mặt phẳng $(Oxy)$.\\
			Do đó $M(1;2;0)$. Suy ra $m=1$, $n=2$, $p=0$.\\
			Vậy $2m-n+2024p=2-2+0=0$.
		\end{itemchoice}
	}
\end{ex}
%Câu 4
\begin{ex}%[2H2V2-2]
	Trong không gian với hệ trục tọa độ $Oxyz$, cho ba điểm $A(1;2;3)$, $B(2;1;5)$, $C(2;4;2)$. 	
	\choiceTF
	{\True Tọa độ trung điểm của $AB$ là $\left(\dfrac{3}{2};\dfrac{3}{2};4\right) $}
	{\True $\overrightarrow{OA}+\overrightarrow{OB}+\overrightarrow{OC}=(5;7;10)$}
	{Góc giữa hai đường thẳng $AB$ và $AC$ bằng $30^\circ$}
	{\True Điểm $I(a;b;c)$ nằm trên mặt phẳng $(Oxz)$ thỏa mãn $T=\left|3\overrightarrow{IB}-\overrightarrow{IC}\right|$ đạt giá trị nhỏ nhất. Khi đó $a-2b+2c=15$}
	\loigiai{
		\begin{itemchoice}
			\itemch Đúng. Gọi $I$ là trung điểm $AB$. Ta có\\
			$\heva{&x_I=\dfrac{x_A+x_B}{2}=\dfrac{1+2}{2}=\dfrac{3}{2}\\&y_I=\dfrac{y_A+y_B}{2}=\dfrac{2+1}{2}=\dfrac{3}{2}\\&z_I=\dfrac{z_A+z_B}{2}=\dfrac{3+5}{2}=4}\Rightarrow I\left(\dfrac{3}{2};\dfrac{3}{2};4\right)$.
			\itemch Đúng. Ta có $\overrightarrow{OA}+\overrightarrow{OB}+\overrightarrow{OC}=(5;7;10)$.
			\itemch Đúng.  Ta có $\overrightarrow{AB}=(1;-1;2)$, $\overrightarrow{AC}=(1;2;-1)$.\\
			\allowdisplaybreaks
			\begin{eqnarray*}
				\cos (AB,AC)&=&\cos\left(\overrightarrow{AB},\overrightarrow{AC}\right)
				\\
				&=&\dfrac{\left|\overrightarrow{AB}\cdot \overrightarrow{AC}\right|}{\left|\overrightarrow{AB}\right|\cdot \left|\overrightarrow{AC}\right|}
				\\
				&=&\dfrac{|1\cdot 1+(-1)\cdot 2+2\cdot (-1)|}{\sqrt{1^2+(-1)^2+2^2}\cdot \sqrt{1^2+2^2+(-1)^2}}
				\\
				&=&\dfrac{1}{2}.
			\end{eqnarray*}
			Suy ra $(AB,AC)=60^\circ$.
			\itemch Sai. Gọi $K(x;y;z)$ thỏa mãn \[3\overrightarrow{KB}-\overrightarrow{KC}=\overrightarrow{0}\Leftrightarrow \heva{&3(2-x)-(2-x)=0\\&3(1-y)-(4-y)=0\\&3(5-z)-(2-z)=0}\Leftrightarrow \heva{&x=2\\&y=-\dfrac{1}{2}\\&z=\dfrac{13}{2}.}\]
			Suy ra $K\left(2;-\dfrac{1}{2};\dfrac{13}{2}\right)$.\\
			Khi đó $T=\left|3\overrightarrow{IB}-\overrightarrow{IC}\right|=\left|3\overrightarrow{IK}+3\overrightarrow{KB}-\overrightarrow{IK}-\overrightarrow{KC}\right|=\left|2\overrightarrow{IK}\right|=2IK$.\\
			$T$ đạt giá trị nhỏ nhất khi và chỉ khi $I$ là hình chiếu của $K$ trên $(Oxz)$.
			Suy ra $ I\left(2;0;\dfrac{13}{2}\right)$.\\
			Suy ra $a=2$, $b=0$, $c=\dfrac{13}{2}$.\\
			 Vậy $a-2b+2c=15$.
		\end{itemchoice}
	}
\end{ex}

\Closesolutionfile{ans}
\indapan{3}{ans/ans-0-B8-DS}
\Opensolutionfile{ans}[ans/ans-0-B8-KQ]
\caukq
\setcounter{ex}{0}
%Câu 1
\begin{ex}%[2H2V2-2]
	Trong không gian $Oxyz$, cho hai điểm $A(-1;-1;1)$ và $B(3;-1;1)$. Điểm $M(a;b;c)$ thỏa mãn $\overrightarrow{AM}=3\overrightarrow{MB}$. Tính $a-2b+c$.
	\shortans{$5$}
	\loigiai{
		Gọi $M(x;y;z)$ ta có $\overrightarrow{AM}=(x+1;y+1;z-1)$ và $\overrightarrow{MB}=(3-x;-1-y;1-z)$.\\
		Do đó
		\[\overrightarrow{AM}=3\overrightarrow{MB} \Leftrightarrow \heva{& x+1=3(3-x) \\ & y+1=3(-1-y) \\ & z-1=3(1-z)}\Leftrightarrow \heva{& x= 2\\ & y=-1 \\ & z=1.}\]
		Suy ra $M(2;-1;1)$. Suy ra $a=2$, $b=-1$, $c=1$ suy ra $a-2b+c=2+2+1=5$.
	}
\end{ex}
%Câu 2
\begin{ex}%[2H2H2-2]
	Trong không gian $Oxyz$, cho ba điểm $A(0;1;3)$, $B(3;2;8)$ và $C(-2 ; m;4)$. Tìm $m$ để tam giác $A B C$ vuông tại $A$.
	\shortans{$2$}
	\loigiai{
		Ta có $\overrightarrow{AB} = (3;1;5)$, $\overrightarrow{AC} = (-2;m-1;1)$. \\
		Tam giác $A B C$ vuông tại $A$ khi và chỉ khi 
		\[
		\overrightarrow{AB} \cdot \overrightarrow{AC} =0 \Leftrightarrow -6+m-1+5 =0 \Leftrightarrow m=2.
		\]
	}
\end{ex}

%Câu 3
\begin{ex}%[2H2V2-2]
	Trong không gian với hệ tọa độ $Oxyz$, cho ba điểm $A(-2 ; 3 ; 1)$, $B\left(\dfrac{1}{4} ; 0 ; 1\right)$, $C(2 ; 0 ; 1)$. Điểm $D(a;b;c)$ là chân đường phân giác trong góc $A$ của tam giác $ABC$. Tính $2022a+2023b-2024c$.
	\shortans{$-2$}
	\loigiai{
		\immini{
			Gọi $D$ là chân đường phân giác trong góc $A$ của tam giác $ABC$, ta có
			\[ \dfrac{BD}{CD}=\dfrac{AB}{AC}=\dfrac{\sqrt{\left(\dfrac{1}{4}+2\right)^2+(0-3)^2+(1-1)^2}}{\sqrt{(2+2)^2+(0-3)^2+(1-1)^2}}= \dfrac{3}{4}.\]
			Từ đó suy ra 
			\[ \overrightarrow{BD}=\dfrac{3}{7}\overrightarrow{BC} \Leftrightarrow \heva{& x_D - \dfrac{1}{4}=\dfrac{3}{7}\left(2-\dfrac{1}{4}\right) \\ 
				& y_D-0 = \dfrac{3}{7}(0-0) \\ & z_D - 1 = \dfrac{3}{7}(1-1)} \Leftrightarrow \heva{& x_D = 1 \\ & y_D = 0 \\ & z_D = 1.}\]
			Suy ra $D(1;0;1)$. Do đó $a=1$, $b=0$, $c=1$.\\
			Vậy $2022a+2023b-2024c=2023+0-2024=-2$.
		}{
			\begin{tikzpicture}[scale=1, font=\footnotesize, line join=round, line cap=round, >=stealth]
				\path (0,0) coordinate (B) (2,3) coordinate (A)  (5,0) coordinate (C)
				coordinate (D) at ([scale around={3/7:(B)}] C);;
				\draw (A)--(B) (A)--(C) (A)--(D) (B)--(C);
				\foreach \p/\r in{A/90,B/-180,C/0,D/-90} \fill (\p) circle (1.2pt) node[shift={(\r:3mm)}]{$\p$};
			\end{tikzpicture}
		}
	}
\end{ex}
%Câu 4
\begin{ex}%[2H2V2-2]
	Trong không gian $Oxyz$, cho tam giác $ABC$ với $A(1;2;5)$, $B(2;4;-3)$, $C(3;3;1)$. Gọi $G$ là trọng tâm của tam giác $ABC$ và $M$ là điểm thay đổi trên mặt phẳng $(Oxy)$. Độ dài $GM$ ngắn nhất bằng bao nhiêu?
	\shortans{$1$}
	\loigiai{
		$G$ là trọng tâm tam giác $ABC$, suy ra $G\left(\dfrac{1+2+3}{3};\dfrac{2+4+3}{3};\dfrac{5+(-3)+1}{3}\right) = \left(2;3;1\right)$.\\
		Gọi $H$ là hình chiếu vuông góc của $G$ trên mặt phẳng $(Oxy)$, khi đó $GH$ là khoảng cách từ $G$ đến mặt phẳng $(Oxy)$, ta có
		$$ GH=\mathrm{d}\left(G,(Oxy)\right)=1.$$
		Với $M$ là điểm thay đổi trên mặt phẳng $(Oxy)$, ta có $GM \geq G H=1$.\\
		Do đó $GM$ ngắn nhất $\Leftrightarrow M \equiv H$.\\
		Vậy độ dài $GM$ ngắn nhất bằng $1$.
	}
\end{ex}
%Câu 5
\begin{ex}%[2H2H2-2]
	Trong không gian với hệ tọa độ $Oxyz$, cho các điểm $E(3;-2;1)$, $F(7;18;-7)$. Điểm $M(a;b;c)$ thỏa mãn $\overrightarrow{FM}=3\overrightarrow{ME}$. Khi đó $a-b+c$ bằng bao nhiêu?
	\shortans{$0$}
	\loigiai{Gọi $M(x;\,y;\,z)$ khi đó $\overrightarrow{FM}= (x-7;\,y-18;\,z+7)$ và $\overrightarrow{ME}=(3-x;\,-2-y;\,1-z)$.\\
		Ta có $\overrightarrow{FM}=3\overrightarrow{ME}\Leftrightarrow (x-7;\,y-18;\,z+7)=(9-3x;\,-6-3y;\,3-3z)$\\
		Từ đó ta có hệ phương trình $\heva{&x-7=9-3x\\&y-18=-6-3y\\&z+7=3-3z}\Leftrightarrow\heva{&x=4\\&y=3\\&z=-1.}$\\
		Suy ra $M(4;\,3;\,-1)$.\\
		Suy ra $a=4$, $b=3$, $c=-1$. Vậy $a-b+c=0$.
	}
\end{ex}
%Câu 6
\begin{ex}%[2H2V2-2]
	Trong không gian $Oxyz$, cho tam giác $ABC$ biết $A(1;2;3)$, $B(0;-1;4)$, $C(2;2;5)$. Điểm $M(x_M;y_M;z_M)$ thuộc mặt phẳng $Oxy$ sao cho biểu thức $S=MA^2+MB^2+MC^2$ đạt giá trị nhỏ nhất. Tính $P=x_M+y_M+z_M$.
	\shortans{$2$}
	\loigiai{
		Gọi $G$ là trọng tâm tam giác $ABC \Rightarrow G(1;1;4)$.\\
		Ta có  
		\allowdisplaybreaks
		$\begin{aligned}[t]
			S=MA^2+MB^2+MC^2
			&=\overrightarrow{MA}^2+\overrightarrow{MB}^2+\overrightarrow{MC}^2\\
			&=\left( \overrightarrow{MG}+\overrightarrow{GA} \right)^2+\left( \overrightarrow{MG}+\overrightarrow{GB} \right)^2+\left( \overrightarrow{MG}+\overrightarrow{GC} \right)^2\\
			&=3\overrightarrow{MG}^2+\overrightarrow{GA}^2+\overrightarrow{GB}^2+\overrightarrow{GC}^2+2\overrightarrow{MG} \cdot \left( \overrightarrow{GA}+\overrightarrow{GB}+\overrightarrow{GC} \right)\\
			&=3MG^2+GA^2+GB^2+GC^2.
		\end{aligned}$\\
		Vì $GA^2+GB^2+GC^2$ không đổi nên $S$ đạt giá trị nhỏ nhất khi và chỉ khi $MG$ nhỏ nhất, tức là $M$ là hình chiếu của $G$ trên mặt phẳng $Oxy$.\\
		Vậy $M(1;1;0)$ và $P=2$.
	}
\end{ex}

\Closesolutionfile{ans}
\indapan{6}{ans/ans-0-B8-KQ}

